\documentclass[12pt,english,a4paper]{report}
\pdfobjcompresslevel=0
\usepackage[usenames,dvipsnames]{xcolor}
\usepackage[includeheadfoot,margin=0.8 in,top=0.6 in]{geometry}
\usepackage{siunitx,physics,cancel,upgreek,varioref,minted,booktabs,tocloft, pdfpages}
\usepackage{mathtools}
\usepackage{babel}
\usepackage{graphicx}
\usepackage{float}
\usepackage{fouriernc}
\usepackage{fancyhdr}
\usepackage[utf8]{inputenc}
\usepackage{amsmath}
\usepackage{amssymb}
\usepackage{textcomp}
\usepackage{lastpage}
\usepackage{microtype}
\usepackage[linktoc=all, bookmarks=true, pdfauthor={Anders Johansson}]{hyperref}
\renewcommand{\CancelColor}{\color{red}}
\renewcommand{\exp}[1]{\mathrm{e}^{#1}}
\newcommand{\tittel}[1]{\title{#1 \vspace{-7ex}}\author{}\date{}\maketitle\thispagestyle{fancy}\pagestyle{fancy}\setcounter{page}{1}}

\newcommand{\deloppg}[2][]{\subsection*{#2) #1}\addcontentsline{toc}{subsection}{#2)}\refstepcounter{subsection}\label{#2}}
\newcommand{\oppg}[1]{\section*{Oppgave #1}\addcontentsline{toc}{section}{Oppgave #1}\refstepcounter{section}\label{oppg#1}}

% \labelformat{section}{seksjon~#1}
% \labelformat{subsection}{avsnitt~#1}
% \labelformat{subsubsection}{avsnitt~#1}
% \labelformat{equation}{likning~(#1)}
% \labelformat{figure}{figur~#1}
% \labelformat{table}{tabell~#1}

\renewcommand{\footrulewidth}{\headrulewidth}
\tocloftpagestyle{fancy}

\setcounter{secnumdepth}{4}
\renewcommand{\thesection}{\arabic{section}}
\renewcommand{\thesubsection}{\arabic{section}.\arabic{subsection}}
\renewcommand{\thesubsubsection}{\arabic{section}.\arabic{subsection}.\arabic{subsubsection}}
\setlength{\parindent}{0cm}
\setlength{\parskip}{1em}

\newcommand{\eqtag}[1]{\refstepcounter{equation}\tag{\theequation}\label{#1}}
\hypersetup{colorlinks=true,urlcolor=blue,linkcolor=black}

\sisetup{detect-all}
\sisetup{exponent-product = \cdot, output-product = \cdot,per-mode=symbol}
\sisetup{output-decimal-marker={.}}
\sisetup{round-mode = off, round-precision=3}
\sisetup{number-unit-product = \ }

\allowdisplaybreaks[4]
\fancyhf{}

\rhead{Anders Johansson}
\rfoot{Page \thepage{} of \pageref{LastPage}}
\lhead{FYS3150}
%
\usepackage[backend=biber,citestyle=numeric-comp,bibstyle=numeric,sorting=none]{biblatex}
\DefineBibliographyStrings{norsk}{%
  bibliography = {Referanser},
}
\addbibresource{kilder.bib}

\begin{document}
%\includepdf{forside.pdf}
\pagestyle{fancy}
\tableofcontents

\section{Abstract}

All files for this project are available at GitHub\footnote{\url{https://github.com/anjohan/Offentlig/tree/master/FYS3150/Oblig2}}.


\section{Introduction}
%%
\section{Obligbesvarelse}
\deloppg{a}
If \(U\) is orthogonal, \(U^TU=I\). This implies:
\[
\vec{w}_i^T\vec{w}_j = \qty(U\vec{v}_i)^T\qty(U\vec{v}_j) = \vec{v}_i^T\overbrace{U^TU}^I\vec{v}_j = \vec{v}_i^T\vec{v}_j=\vec{v}_i\cdot \vec{v}_j
\]
Orthogonal transformations hence conserve the inner product, and therefore also the orthogonality.

\begin{figure}[H]
\centering
% GNUPLOT: LaTeX picture with Postscript
\begingroup
  \makeatletter
  \providecommand\color[2][]{%
    \GenericError{(gnuplot) \space\space\space\@spaces}{%
      Package color not loaded in conjunction with
      terminal option `colourtext'%
    }{See the gnuplot documentation for explanation.%
    }{Either use 'blacktext' in gnuplot or load the package
      color.sty in LaTeX.}%
    \renewcommand\color[2][]{}%
  }%
  \providecommand\includegraphics[2][]{%
    \GenericError{(gnuplot) \space\space\space\@spaces}{%
      Package graphicx or graphics not loaded%
    }{See the gnuplot documentation for explanation.%
    }{The gnuplot epslatex terminal needs graphicx.sty or graphics.sty.}%
    \renewcommand\includegraphics[2][]{}%
  }%
  \providecommand\rotatebox[2]{#2}%
  \@ifundefined{ifGPcolor}{%
    \newif\ifGPcolor
    \GPcolortrue
  }{}%
  \@ifundefined{ifGPblacktext}{%
    \newif\ifGPblacktext
    \GPblacktextfalse
  }{}%
  % define a \g@addto@macro without @ in the name:
  \let\gplgaddtomacro\g@addto@macro
  % define empty templates for all commands taking text:
  \gdef\gplbacktext{}%
  \gdef\gplfronttext{}%
  \makeatother
  \ifGPblacktext
    % no textcolor at all
    \def\colorrgb#1{}%
    \def\colorgray#1{}%
  \else
    % gray or color?
    \ifGPcolor
      \def\colorrgb#1{\color[rgb]{#1}}%
      \def\colorgray#1{\color[gray]{#1}}%
      \expandafter\def\csname LTw\endcsname{\color{white}}%
      \expandafter\def\csname LTb\endcsname{\color{black}}%
      \expandafter\def\csname LTa\endcsname{\color{black}}%
      \expandafter\def\csname LT0\endcsname{\color[rgb]{1,0,0}}%
      \expandafter\def\csname LT1\endcsname{\color[rgb]{0,1,0}}%
      \expandafter\def\csname LT2\endcsname{\color[rgb]{0,0,1}}%
      \expandafter\def\csname LT3\endcsname{\color[rgb]{1,0,1}}%
      \expandafter\def\csname LT4\endcsname{\color[rgb]{0,1,1}}%
      \expandafter\def\csname LT5\endcsname{\color[rgb]{1,1,0}}%
      \expandafter\def\csname LT6\endcsname{\color[rgb]{0,0,0}}%
      \expandafter\def\csname LT7\endcsname{\color[rgb]{1,0.3,0}}%
      \expandafter\def\csname LT8\endcsname{\color[rgb]{0.5,0.5,0.5}}%
    \else
      % gray
      \def\colorrgb#1{\color{black}}%
      \def\colorgray#1{\color[gray]{#1}}%
      \expandafter\def\csname LTw\endcsname{\color{white}}%
      \expandafter\def\csname LTb\endcsname{\color{black}}%
      \expandafter\def\csname LTa\endcsname{\color{black}}%
      \expandafter\def\csname LT0\endcsname{\color{black}}%
      \expandafter\def\csname LT1\endcsname{\color{black}}%
      \expandafter\def\csname LT2\endcsname{\color{black}}%
      \expandafter\def\csname LT3\endcsname{\color{black}}%
      \expandafter\def\csname LT4\endcsname{\color{black}}%
      \expandafter\def\csname LT5\endcsname{\color{black}}%
      \expandafter\def\csname LT6\endcsname{\color{black}}%
      \expandafter\def\csname LT7\endcsname{\color{black}}%
      \expandafter\def\csname LT8\endcsname{\color{black}}%
    \fi
  \fi
    \setlength{\unitlength}{0.0500bp}%
    \ifx\gptboxheight\undefined%
      \newlength{\gptboxheight}%
      \newlength{\gptboxwidth}%
      \newsavebox{\gptboxtext}%
    \fi%
    \setlength{\fboxrule}{0.5pt}%
    \setlength{\fboxsep}{1pt}%
\begin{picture}(7200.00,4320.00)%
    \gplgaddtomacro\gplbacktext{%
      \csname LTb\endcsname%
      \put(1078,704){\makebox(0,0)[r]{\strut{}$0$}}%
      \csname LTb\endcsname%
      \put(1078,1263){\makebox(0,0)[r]{\strut{}$0.005$}}%
      \csname LTb\endcsname%
      \put(1078,1821){\makebox(0,0)[r]{\strut{}$0.01$}}%
      \csname LTb\endcsname%
      \put(1078,2380){\makebox(0,0)[r]{\strut{}$0.015$}}%
      \csname LTb\endcsname%
      \put(1078,2938){\makebox(0,0)[r]{\strut{}$0.02$}}%
      \csname LTb\endcsname%
      \put(1078,3497){\makebox(0,0)[r]{\strut{}$0.025$}}%
      \csname LTb\endcsname%
      \put(1078,4055){\makebox(0,0)[r]{\strut{}$0.03$}}%
      \csname LTb\endcsname%
      \put(1210,484){\makebox(0,0){\strut{}$0$}}%
      \csname LTb\endcsname%
      \put(1909,484){\makebox(0,0){\strut{}$1$}}%
      \csname LTb\endcsname%
      \put(2608,484){\makebox(0,0){\strut{}$2$}}%
      \csname LTb\endcsname%
      \put(3307,484){\makebox(0,0){\strut{}$3$}}%
      \csname LTb\endcsname%
      \put(4007,484){\makebox(0,0){\strut{}$4$}}%
      \csname LTb\endcsname%
      \put(4706,484){\makebox(0,0){\strut{}$5$}}%
      \csname LTb\endcsname%
      \put(5405,484){\makebox(0,0){\strut{}$6$}}%
      \csname LTb\endcsname%
      \put(6104,484){\makebox(0,0){\strut{}$7$}}%
      \csname LTb\endcsname%
      \put(6803,484){\makebox(0,0){\strut{}$8$}}%
    }%
    \gplgaddtomacro\gplfronttext{%
      \csname LTb\endcsname%
      \put(176,2379){\rotatebox{-270}{\makebox(0,0){\strut{}$u(\rho)^2$}}}%
      \put(4006,154){\makebox(0,0){\strut{}$\rho$}}%
      \csname LTb\endcsname%
      \put(5816,3882){\makebox(0,0)[r]{\strut{}$\lambda=3.000$}}%
      \csname LTb\endcsname%
      \put(5816,3662){\makebox(0,0)[r]{\strut{}$\lambda=6.998$}}%
      \csname LTb\endcsname%
      \put(5816,3442){\makebox(0,0)[r]{\strut{}$\lambda=10.996$}}%
    }%
    \gplbacktext
    \put(0,0){\includegraphics{plot1}}%
    \gplfronttext
  \end{picture}%
\endgroup

\caption{Simulation of one particle in a harmonic oscillator potential for \(n=250\) with \(\rho_{\max}=8\).}
\end{figure}

\begin{table}[H]
\[
\input{egenverdier.dat}
\]
\caption{The three lowest eigenvalues found by the algorithm for the different number of mesh points \(n\). The analytical values are \(\lambda_0=3\), \(\lambda_1=7\) and \(\lambda_2=11\).}
\end{table}
The number of iterations is approximately quadrupled when the number of mesh points is doubled, indicating that the numbers of iterations runs as \(n^2\). This means that the number of iterations is proportional to the number of elements in the matrix. Note that the matrices handled in this project are tridiagonal, and the trends may be different for dense matrices.


\begin{figure}[H]
\centering
% GNUPLOT: LaTeX picture with Postscript
\begingroup
  \makeatletter
  \providecommand\color[2][]{%
    \GenericError{(gnuplot) \space\space\space\@spaces}{%
      Package color not loaded in conjunction with
      terminal option `colourtext'%
    }{See the gnuplot documentation for explanation.%
    }{Either use 'blacktext' in gnuplot or load the package
      color.sty in LaTeX.}%
    \renewcommand\color[2][]{}%
  }%
  \providecommand\includegraphics[2][]{%
    \GenericError{(gnuplot) \space\space\space\@spaces}{%
      Package graphicx or graphics not loaded%
    }{See the gnuplot documentation for explanation.%
    }{The gnuplot epslatex terminal needs graphicx.sty or graphics.sty.}%
    \renewcommand\includegraphics[2][]{}%
  }%
  \providecommand\rotatebox[2]{#2}%
  \@ifundefined{ifGPcolor}{%
    \newif\ifGPcolor
    \GPcolortrue
  }{}%
  \@ifundefined{ifGPblacktext}{%
    \newif\ifGPblacktext
    \GPblacktextfalse
  }{}%
  % define a \g@addto@macro without @ in the name:
  \let\gplgaddtomacro\g@addto@macro
  % define empty templates for all commands taking text:
  \gdef\gplbacktext{}%
  \gdef\gplfronttext{}%
  \makeatother
  \ifGPblacktext
    % no textcolor at all
    \def\colorrgb#1{}%
    \def\colorgray#1{}%
  \else
    % gray or color?
    \ifGPcolor
      \def\colorrgb#1{\color[rgb]{#1}}%
      \def\colorgray#1{\color[gray]{#1}}%
      \expandafter\def\csname LTw\endcsname{\color{white}}%
      \expandafter\def\csname LTb\endcsname{\color{black}}%
      \expandafter\def\csname LTa\endcsname{\color{black}}%
      \expandafter\def\csname LT0\endcsname{\color[rgb]{1,0,0}}%
      \expandafter\def\csname LT1\endcsname{\color[rgb]{0,1,0}}%
      \expandafter\def\csname LT2\endcsname{\color[rgb]{0,0,1}}%
      \expandafter\def\csname LT3\endcsname{\color[rgb]{1,0,1}}%
      \expandafter\def\csname LT4\endcsname{\color[rgb]{0,1,1}}%
      \expandafter\def\csname LT5\endcsname{\color[rgb]{1,1,0}}%
      \expandafter\def\csname LT6\endcsname{\color[rgb]{0,0,0}}%
      \expandafter\def\csname LT7\endcsname{\color[rgb]{1,0.3,0}}%
      \expandafter\def\csname LT8\endcsname{\color[rgb]{0.5,0.5,0.5}}%
    \else
      % gray
      \def\colorrgb#1{\color{black}}%
      \def\colorgray#1{\color[gray]{#1}}%
      \expandafter\def\csname LTw\endcsname{\color{white}}%
      \expandafter\def\csname LTb\endcsname{\color{black}}%
      \expandafter\def\csname LTa\endcsname{\color{black}}%
      \expandafter\def\csname LT0\endcsname{\color{black}}%
      \expandafter\def\csname LT1\endcsname{\color{black}}%
      \expandafter\def\csname LT2\endcsname{\color{black}}%
      \expandafter\def\csname LT3\endcsname{\color{black}}%
      \expandafter\def\csname LT4\endcsname{\color{black}}%
      \expandafter\def\csname LT5\endcsname{\color{black}}%
      \expandafter\def\csname LT6\endcsname{\color{black}}%
      \expandafter\def\csname LT7\endcsname{\color{black}}%
      \expandafter\def\csname LT8\endcsname{\color{black}}%
    \fi
  \fi
    \setlength{\unitlength}{0.0500bp}%
    \ifx\gptboxheight\undefined%
      \newlength{\gptboxheight}%
      \newlength{\gptboxwidth}%
      \newsavebox{\gptboxtext}%
    \fi%
    \setlength{\fboxrule}{0.5pt}%
    \setlength{\fboxsep}{1pt}%
\begin{picture}(7200.00,4320.00)%
    \gplgaddtomacro\gplbacktext{%
      \csname LTb\endcsname%
      \put(1078,704){\makebox(0,0)[r]{\strut{}$0$}}%
      \csname LTb\endcsname%
      \put(1078,1263){\makebox(0,0)[r]{\strut{}$0.005$}}%
      \csname LTb\endcsname%
      \put(1078,1821){\makebox(0,0)[r]{\strut{}$0.01$}}%
      \csname LTb\endcsname%
      \put(1078,2380){\makebox(0,0)[r]{\strut{}$0.015$}}%
      \csname LTb\endcsname%
      \put(1078,2938){\makebox(0,0)[r]{\strut{}$0.02$}}%
      \csname LTb\endcsname%
      \put(1078,3497){\makebox(0,0)[r]{\strut{}$0.025$}}%
      \csname LTb\endcsname%
      \put(1078,4055){\makebox(0,0)[r]{\strut{}$0.03$}}%
      \csname LTb\endcsname%
      \put(1210,484){\makebox(0,0){\strut{}$0$}}%
      \csname LTb\endcsname%
      \put(1909,484){\makebox(0,0){\strut{}$1$}}%
      \csname LTb\endcsname%
      \put(2608,484){\makebox(0,0){\strut{}$2$}}%
      \csname LTb\endcsname%
      \put(3307,484){\makebox(0,0){\strut{}$3$}}%
      \csname LTb\endcsname%
      \put(4007,484){\makebox(0,0){\strut{}$4$}}%
      \csname LTb\endcsname%
      \put(4706,484){\makebox(0,0){\strut{}$5$}}%
      \csname LTb\endcsname%
      \put(5405,484){\makebox(0,0){\strut{}$6$}}%
      \csname LTb\endcsname%
      \put(6104,484){\makebox(0,0){\strut{}$7$}}%
      \csname LTb\endcsname%
      \put(6803,484){\makebox(0,0){\strut{}$8$}}%
    }%
    \gplgaddtomacro\gplfronttext{%
      \csname LTb\endcsname%
      \put(176,2379){\rotatebox{-270}{\makebox(0,0){\strut{}$u(\rho)^2$}}}%
      \put(4006,154){\makebox(0,0){\strut{}$\rho$}}%
      \csname LTb\endcsname%
      \put(5816,3882){\makebox(0,0)[r]{\strut{}Uten interaksjon}}%
      \csname LTb\endcsname%
      \put(5816,3662){\makebox(0,0)[r]{\strut{}Med interaksjon}}%
    }%
    \gplbacktext
    \put(0,0){\includegraphics{plot2}}%
    \gplfronttext
  \end{picture}%
\endgroup

\caption{Simulation of two particles, with and without Coulomb interaction, for two different values of \(\omega\). A larger \(\omega\) corresponds to a stronger harmonic oscillator.}
\end{figure}











\clearpage
\addcontentsline{toc}{section}{Referanser}
\printbibliography


\end{document}
